\section{When did discrete representation get so bad?}

Flow-based policies have become a prominent choice in robot learning. Discrete
autoregressive alternatives, meanwhile, face a familiar objection: turning
continuous actions into tokens must sacrifice precision.

Yet discrete representations support capable generative models in other
modalities, including audio. Discreteness alone is an incomplete explanation for
why an action policy succeeds or fails.

The useful question is narrower: is the reconstruction error introduced by a
tokenizer large enough to explain the policy’s failures? A tokenizer can
reconstruct demonstrations accurately while producing sequences that a policy
struggles to predict. Conversely, a small reconstruction error can matter a
great deal near contact or a task boundary.

\begin{quote}
Reconstruction quality and generation difficulty are different properties of a
representation.
\end{quote}

Our draft’s central hypothesis is that uneven information allocation and
accumulated prediction errors can explain failures that reconstruction metrics
miss. Establishing that claim requires comparing decoded expert actions with
actions generated by the policy, in the same tasks and under matched conditions.
